% ============================================================================
% Reading guide --- Technical Reference
% Page target: ~2 pages
% ============================================================================

\section*{Reading guide}
\addcontentsline{toc}{section}{Reading guide}

The document is structured in four parts (Foundations, Reference
Architecture on Open Foundations, Migration Playbook for Mixed
Estates, References and Operational Material) and is designed to
be read in order. Each chapter is, however, written so that it
remains intelligible on its own; cross-references are explicit so
the reader can jump in wherever the table of audience entry
points below directs.

\paragraph{Document structure at a glance.}
\begin{itemize}[leftmargin=1.5em,nosep]
  \item \textbf{Part~I --- Foundations} (Ch.~2--4): conceptual frame,
    the sovereignty grid and its twenty indicators, and the
    self-assessment instrument.
  \item \textbf{Part~II --- Reference architecture} (Ch.~5--16):
    architectural principles, the nine layer chapters (identity,
    policy, network, compute, storage, endpoint, observability,
    cryptographic services, orchestration), the sizing chapter and
    the conformance suite.
  \item \textbf{Part~III --- Migration playbook} (Ch.~17--30):
    reading the playbook, the ten waves (Wave~0--Wave~9),
    cross-cutting topics, the reversibility playbook, and the
    coexistence-pattern library.
  \item \textbf{Part~IV --- References and operational material}
    (Ch.~31--33): institutional precedents, external references and
    standards, and the regulatory mapping.
\end{itemize}

\paragraph{Audiences and entry points.} The table below proposes a
reading path for each principal audience. The synthesis chapter
(Chapter~1) is the recommended common starting point.

\begin{table}[H]
\centering\small
\renewcommand{\arraystretch}{1.18}
\begin{tabularx}{\textwidth}{%
  >{\hsize=0.55\hsize\linewidth=\hsize\bfseries\raggedright\arraybackslash}X%
  >{\hsize=1.45\hsize\linewidth=\hsize\raggedright\arraybackslash}X%
}
\toprule
Audience & Recommended reading path \\
\midrule
Governing board / executive leadership
  & Ch.~1 (synthesis), Ch.~3 (sovereignty grid), Ch.~4
    (self-assessment), Ch.~33 (regulatory mapping). \\
CIO / CTO
  & Ch.~1, Ch.~2 (TAC), Ch.~3, Ch.~5 (architectural principles),
    Ch.~17 (playbook reading), Ch.~31 (institutional precedents),
    Ch.~33. \\
Chief Information Security Officer
  & Ch.~1, Ch.~6 (Identity), Ch.~7 (Policy), Ch.~12 (Observability),
    Ch.~13 (Cryptographic services), Ch.~25 (Wave~7 audit and IR),
    Ch.~33. \\
Network architect
  & Ch.~1, Ch.~5, Ch.~8 (Network), Ch.~21 (Wave~3 network),
    Ch.~30 (Coexistence patterns). \\
Platform / SRE engineer
  & Ch.~1, Ch.~9--11 (Compute, Storage, Endpoint),
    Ch.~14 (Orchestration), Ch.~15 (Sizing),
    Ch.~16 (Conformance). \\
Data-protection officer
  & Ch.~1, Ch.~3, Ch.~10 (Storage), Ch.~12,
    Ch.~33 (regulatory mapping: data-protection, cross-border-transfer,
    and national-transposition sections). \\
Programme / migration lead
  & Ch.~1, Ch.~17, Part~III in full,
    Ch.~28--30 (Cross-cutting, Reversibility, Coexistence). \\
Phase-2 peer institution
  & Ch.~1, Ch.~31 (precedents), Ch.~32 (external references),
    Ch.~33 (regulatory mapping), then the chapters most relevant
    to the institution's prevailing brownfield. \\
\bottomrule
\end{tabularx}
\end{table}

\paragraph{Reading the document linearly.} The four parts are
designed to be read in order. Part~I establishes the conceptual
vocabulary on which Parts~II and~III rest; Part~II defines the
architecture against which Part~III's waves measure progress; and
Part~IV consolidates the references and the regulatory anchoring
that the rest of the document points at. A reader new to the
material is best served by following this order.

\paragraph{Reading the document in the book context.} When this
material is encountered as Part~IV of the book \emph{Complete
Dossier}, the reader sees only the synthesis chapter (Chapter~1)
together with an annotated table of contents pointing back to the
present standalone document. The standalone Technical Reference
is the canonical full version: the synthesis chapter is identical
in both contexts, and all cross-references named ``\emph{Chapter
X of the Technical Reference}'' resolve into this document.

\paragraph{Tonality and box conventions.} Coloured boxes have
stable semantics throughout the document:
\begin{itemize}[leftmargin=1.5em,nosep]
  \item \textbf{Blue (Interface contract)} --- normative. The
    minimum set of properties a layer must expose for the
    architecture to compose.
  \item \textbf{Green (Reference instantiation)} --- illustrative.
    The open-foundation default shown sufficient to satisfy the
    contract.
  \item \textbf{Slate (Alternative)} --- illustrative. A
    commercial or hybrid option evaluated against the contract.
  \item \textbf{Red (Common pitfall)} --- cautionary. Patterns
    that have produced poor outcomes and the mitigations that
    address them.
  \item \textbf{Plum (Real-world reference)} --- precedent. A
    documented institution or pattern; statements are framed
    conditionally pending verification by member institutions.
  \item \textbf{Slate caveatbox} --- epistemic posture. A note
    on what the document claims and what it does not.
\end{itemize}

\paragraph{Voice conventions.} The Reference is assertive when
discussing contracts, standards and measurable properties; it is
analytical, not advocatory, when comparing instantiations; and
it is conditional (``would'', ``may'', ``is expected to'') when
projecting onto institutional outcomes that cannot be
demonstrated from the state of the art at the time of writing.
This convention is consistent across the document and is part of
the always-living posture stated in the foreword.

\newpage
